%%%%%%%%%%%%%%%%%%%%%%%%%%%%%%%%%
%
%       EXERCÍCIO
%
%%%%%%%%%%%%%%%%%%%%%%%%%%%%%%%%%

\definevalorquestao

\begin{exercicioBanco}[\valorquestao]
    Para famílias de um certo bairro de Contagem, apresentamos abaixo a tabela de frequência conjunta das variáveis: \textit{número de automóveis} (A) e \textit{número de televisões} (T)
    %
    \[
    \begin{array}{|c|c|c|c|c|}
        \hline
        A\backslash T & 0 & 1 & 2 & \text{Total} \\
        \hline
        0 & 110 & 235 & 120 & 465 \\
        \hline
        1 & 51 & 122 & 178 & 351 \\
        \hline
        2 & 15 & 84 & 162 & 261 \\
        \hline
        \text{Total} & 176 & 441 & 460 & 1077 \\
        \hline
    \end{array}
    \]
    %
    Determine as tabelas marginas de \(A\) e \(T\) e calcule suas médias.
\end{exercicioBanco}

